\section{Aim \& Clinical Motivation}
\label{sec:aim}

Parkinson's disease (PD) is a progressive disorder of the brain's movement-control
circuitry. As the dopamine-producing neurons of the \emph{substantia nigra} die off,
the messages that normally let movement flow smoothly become noisy and unreliable.
The familiar outward signs --- a resting tremor, stiffness, and a general slowing of
movement --- are only part of the picture. For many people in the middle and later
stages of the disease, the most frightening symptom is one that the textbook
descriptions barely capture: their feet simply stop.

\subsection{Freeze of gait: feet glued to the floor}

\emph{Freeze of gait} (FoG) is a sudden, brief inability to take the next step. People
who live with it describe the sensation vividly --- as though their feet have been
\emph{glued to the floor} --- even while the upper body still wants to move forward.
The intention to walk is intact; the legs will not answer. Episodes are typically
short, lasting only seconds, and they are often triggered by exactly the moments where
a person most needs their feet: turning around, stepping through a narrow doorway,
starting to walk from standing, or approaching a chair.

A crucial and counter-intuitive detail is that a freeze is \emph{not} the same as
standing still. During a freeze the feet are usually trembling and shuffling rapidly
\emph{in place} --- a fast, small-amplitude tremor of roughly \SI{3}{}--\SI{8}{\hertz}
--- while the body fails to make forward progress. This rapid trembling-in-place is
the tell-tale mechanical signature that an instrument can latch onto, and it is what
distinguishes a genuine freeze from someone calmly waiting at a kerb. It is also the
single most important fact behind every design decision in this report.

\begin{intuition}
A freeze is not a pause. When you stand at rest, your feet are quiet. During a freeze
the feet are \emph{busy} --- buzzing in place at \SI{3}{}--\SI{8}{\hertz} --- yet you
go nowhere. ``Quiet feet, no progress'' is rest; ``busy feet, no progress'' is a
freeze. That contrast is exactly what an ankle accelerometer can hear.
\end{intuition}

\paragraph{Why it is dangerous.} Freezing is one of the leading causes of falls in
Parkinson's disease. When the feet stop but the body's momentum does not, the person
pitches forward with nothing to catch them. The consequences cascade far beyond the
fall itself. A single fracture can mean surgery, a long recovery, and a lasting loss
of confidence; the fear of the \emph{next} fall then shrinks a person's world. Many
people respond by walking less, going out less, and leaning more heavily on carers ---
so freezing quietly erodes the very independence that mobility represents. Because
episodes are brief and unpredictable, they are also hard to capture in a clinic: a
patient may walk perfectly down a corridor for a neurologist and freeze repeatedly at
home an hour later. What the clinician needs, and rarely gets, is an honest record of
\emph{when} and \emph{how often} freezing actually happens during ordinary life.

\begin{keyidea}
Freezing is brief, unpredictable, and hard to reproduce on demand --- which is exactly
why a short clinic visit underestimates it. The clinical value of a wearable is not
that it does anything clever in the moment; it is that it watches \emph{continuously}
and turns scattered, invisible episodes into a record a clinician can act on.
\end{keyidea}

\subsection{Why a cheap, continuous wearable matters}

If freezing only mattered in the consulting room, an expensive lab gait mat would be
enough. But freezing happens at home, in the supermarket, in the doorway at night ---
precisely where no instrument is watching. A wearable changes that arithmetic. It is
worn all day, it costs little, and it converts a symptom that is otherwise reported
from fading memory into objective, time-stamped data. A low-cost device can be given to
many patients at once; a continuous one captures the rare episodes that a snapshot
misses. Together, ``cheap'' and ``continuous'' are what make population-scale, real-life
freeze monitoring realistic rather than aspirational.

\subsection{Cadence: the aim, stated precisely}

\textbf{Cadence} is a low-cost, ankle-worn wearable whose aim is narrow and deliberate:
to \emph{faithfully detect} freeze-of-gait episodes from a \emph{single ankle
accelerometer}, and to \emph{intuitively display} them --- on the device itself and on a
live dashboard --- so that the person and their clinician can see when freezing happened.
It is built on an Adafruit Circuit Playground Express, worn in a soft ankle sleeve, and
it does its sensing and decision-making from the three axes of one tiny accelerometer.
Two verbs carry the whole project: \emph{detect} and \emph{display}.

\begin{takeaway}
Cadence has exactly one job: to \emph{watch} the ankle, recognise the
\SI{3}{}--\SI{8}{\hertz} signature of a freeze, and \emph{show} it clearly. Detect, and
display --- nothing more.
\end{takeaway}

\subsection{A monitor, not a therapy}

This positioning must be crisp, because it is easy to get wrong. There is a well-known
class of Parkinson's aids called \emph{cueing therapies}: devices that, when they sense
a freeze, fire a stimulus --- a rhythmic buzz, a beep, a flash --- to help the person
break out of it and step again. Cueing is an \emph{intervention}: it acts \emph{on} the
patient in the moment.

\textbf{Cadence is not one of these.} It is a \emph{monitor}. Its purpose is to observe
and record freezing faithfully, not to interrupt or treat it. A coin vibration motor is
physically fitted to the board, but it is \emph{deliberately disabled in firmware}: the
hardware exists, but no firmware path ever drives it. When Cadence recognises a freeze it
\emph{reports} that fact --- on an OLED screen, on the onboard NeoPixel LED ring, and on
a live web dashboard --- and it does \emph{not} buzz, cue, or otherwise act on the wearer.
The distinction between observing and intervening is the line Cadence is careful never to
cross. Table~\ref{tab:monitor-vs-cue} makes the contrast explicit.

\begin{table}[H]
\centering
\caption{Cadence is a monitor, not a cueing therapy. The two aim at different goals and
are evaluated by different yardsticks.}
\label{tab:monitor-vs-cue}
\small
\begin{tabular}{@{}p{0.30\linewidth} p{0.31\linewidth} p{0.31\linewidth}@{}}
\toprule
 & \textbf{Cueing therapy} & \textbf{Cadence (monitor)} \\
\midrule
Goal            & Help end a freeze \emph{now} & Record \emph{when} freezes happen \\
Acts on patient & Yes --- fires a cue          & No --- only reports \\
On a freeze     & Buzz / beep / flash to cue   & Show on OLED + LEDs + dashboard \\
Coin motor      & Driven                       & Fitted but \emph{disabled} \\
Success means    & Patient steps again          & Faithful, honest detection \\
\bottomrule
\end{tabular}
\end{table}

\subsection{What a freeze looks like in the data}

Before diving into signal processing in later sections, it helps to see, once, what
Cadence is actually looking at. Figure~\ref{fig:gaitstrip} sketches the canonical
sequence Cadence is built to follow: a calm \emph{still} period, ordinary \emph{walking},
a sudden \emph{freeze}, and a return to walking. The accelerometer trace tells the story
in amplitude alone --- the freeze is the burst of fast, jagged activity sandwiched
between two stretches of smoother motion.

\begin{figure}[H]
\centering
\resizebox{0.96\textwidth}{!}{%
\begin{tikzpicture}[font=\sffamily]

% ---- geometry ----
\def\H{2.6}      % strip height
\def\W{15.2}     % total width
% phase boundaries (x positions)
\def\xa{0}       % start
\def\xb{3.6}     % still -> walk
\def\xc{7.6}     % walk -> freeze
\def\xd{10.4}    % freeze -> walk
\def\xe{15.2}    % end

% ---- coloured phase bands ----
\fill[palenavy]   (\xa,0) rectangle (\xb,\H);
\fill[paleteal]   (\xb,0) rectangle (\xc,\H);
\fill[palefreeze] (\xc,0) rectangle (\xd,\H);
\fill[paleteal]   (\xd,0) rectangle (\xe,\H);
\draw[rule,line width=0.4pt] (\xa,0) rectangle (\xe,\H);

% baseline (mean) for the trace
\def\mid{1.3}

% ---- STILL: tiny quiet noise ----
\draw[ink,line width=0.9pt]
  (\xa,\mid)
  (0.20,1.34) -- (0.55,1.27) -- (0.90,1.33) -- (1.25,1.28) -- (1.60,1.32)
  -- (1.95,1.27) -- (2.30,1.33) -- (2.65,1.28) -- (3.00,1.32) -- (3.35,1.29)
  -- (\xb,1.30);

% ---- WALK: smooth, moderate oscillation (heel strikes) ----
\draw[grn,line width=1.0pt]
  (\xb,\mid)
  -- (3.95,1.78) -- (4.30,0.86) -- (4.65,1.80) -- (5.00,0.84)
  -- (5.35,1.79) -- (5.70,0.85) -- (6.05,1.80) -- (6.40,0.84)
  -- (6.75,1.79) -- (7.10,0.86) -- (7.45,1.55) -- (\xc,1.30);

% ---- FREEZE: fast, large, jagged trembling in place ----
% (upper peaks held below ~2.05 so the FREEZE label above stays legible;
%  still clearly the tallest, most jagged trace — walking peaks at 1.80)
\draw[freeze,line width=1.0pt]
  (\xc,\mid)
  -- (7.78,1.98) -- (7.96,0.45) -- (8.14,2.04) -- (8.32,0.38)
  -- (8.50,2.06) -- (8.68,0.35) -- (8.86,2.06) -- (9.04,0.36)
  -- (9.22,2.05) -- (9.40,0.38) -- (9.58,2.04) -- (9.76,0.42)
  -- (9.94,2.00) -- (10.12,0.55) -- (\xd,1.30);

% ---- WALK again ----
\draw[grn,line width=1.0pt]
  (\xd,\mid)
  -- (10.78,1.78) -- (11.13,0.86) -- (11.48,1.80) -- (11.83,0.84)
  -- (12.18,1.79) -- (12.53,0.85) -- (12.88,1.80) -- (13.23,0.84)
  -- (13.58,1.79) -- (13.93,0.85) -- (14.28,1.80) -- (14.63,0.86)
  -- (14.98,1.55) -- (\xe,1.30);

% ---- phase labels (top) ----
\node[anchor=north,navy,font=\bfseries] at ({(\xa+\xb)/2},\H-0.08) {STILL};
\node[anchor=north,teal,font=\bfseries] at ({(\xb+\xc)/2},\H-0.08) {WALKING};
\node[anchor=north,freeze,font=\bfseries] at ({(\xc+\xd)/2},\H-0.08) {FREEZE};
\node[anchor=north,teal,font=\bfseries] at ({(\xd+\xe)/2},\H-0.08) {WALKING};

% ---- descriptive sub-labels (bottom, inside) ----
\node[anchor=south,mute,font=\scriptsize] at ({(\xa+\xb)/2},0.10) {quiet feet};
\node[anchor=south,mute,font=\scriptsize] at ({(\xb+\xc)/2},0.10) {steady heel strikes};
\node[anchor=south,freeze,font=\scriptsize\bfseries] at ({(\xc+\xd)/2},0.10) {trembling in place 3--8\,Hz};
\node[anchor=south,mute,font=\scriptsize] at ({(\xd+\xe)/2},0.10) {steady heel strikes};

% ---- y-axis hint ----
\node[anchor=east,mute,font=\scriptsize,rotate=90] at (-0.30,\H/2)
  {ankle $\norm{a}$};

% ---- time arrow under the strip ----
\draw[-{Latex[length=3mm]},mute,line width=0.8pt]
  (\xa,-0.55) -- (\xe,-0.55);
\node[anchor=west,mute,font=\scriptsize] at (\xe,-0.55) {time};

% ---- annotation bracket over the freeze ----
\draw[decorate,decoration={brace,amplitude=5pt,mirror},freeze,line width=0.7pt]
  (\xc,-0.18) -- (\xd,-0.18);
\node[anchor=north,freeze,font=\scriptsize\itshape] at ({(\xc+\xd)/2},-0.30)
  {what Cadence must catch};

\end{tikzpicture}}
\caption{The canonical gait sequence Cadence is built to follow.\ Calm \emph{still}
periods and ordinary \emph{walking} both produce orderly traces; a \emph{freeze} appears
as a burst of fast, large-amplitude trembling ``in place'' before walking resumes.\ The
red episode --- not the quiet stretches --- is what the device must detect.}
\label{fig:gaitstrip}
\end{figure}

That schematic is not just an illustration; it matches what the hardware actually
records. Figure~\ref{fig:fiperphase} shows the central decision quantity --- the
\emph{Freeze Index}, the ratio of fast (\SI{3}{}--\SI{8}{\hertz}) to slow
(\SI{0.5}{}--\SI{3}{\hertz}) energy in the ankle signal --- measured on a real worksheet
capture that runs through exactly this still\,$\rightarrow$\,walk\,$\rightarrow$\,freeze\,$\rightarrow$\,walk
sequence. Standing still and ordinary walking both sit around a Freeze Index of roughly
$1.1$ and $1.2$; the freeze phase leaps to a mean of $6.89$, far above the deployed
trip line of $1.815$. The full mechanics of how that index is computed, thresholded, and
debounced are the subject of later sections; here it is enough to see that the freeze
genuinely stands apart from everything around it.

\begin{figure}[H]
\centering
\includegraphics[width=0.82\linewidth]{fi_per_phase}
\caption{Real Freeze Index per phase on the worksheet capture (\code{capture1}, \SI{64}{\hertz},
\SI{41.4}{\second}).\ Still and walking phases stay near $1.1$--$1.2$; the freeze phase
jumps to a mean of $6.89$, well clear of the deployed threshold of $1.815$ (dashed).\ The
ordered gait sequence of Figure~\ref{fig:gaitstrip} is visible directly in the data.}
\label{fig:fiperphase}
\end{figure}

\begin{keyidea}
The freeze does not hide. On real data the Freeze Index of a freeze ($\approx 6.9$) is
several times that of still or walking ($\approx 1.1$--$1.2$). Cadence's job is to draw a
sensible line between them and to report --- never to interfere.
\end{keyidea}

\subsection{What this report covers}

The remainder of this report follows Cadence from physical signal to clinical readout.

\begin{itemize}
\item \textbf{Sensing.} The accelerometer, why it is sampled at \SI{64}{\hertz} over
$\pm\SI{8}{\g}$, and how three axes are fused into one orientation-invariant magnitude.
\item \textbf{Signal processing.} Welch power spectra, the locomotor / freeze /
rest-tremor frequency bands, and the Freeze Index that turns a wiggling trace into a
single number.
\item \textbf{The shipped detector.} The transparent ``glass-box'' rule --- an energy
gate, a Freeze-Index threshold, and a debounce filter --- and how its operating point was
chosen to favour catching freezes.
\item \textbf{A learned alternative.} \emph{FoGNet}, a small 1-D convolutional network
explored as a ``given more time'' upgrade, with its honest cross-validated results.
\item \textbf{Evaluation.} Why sensitivity and specificity, not raw accuracy, are the
metrics that matter when freezes are rare.
\item \textbf{Hardware \& software.} The Circuit Playground build, the live dashboard,
and the firmware that records and displays --- but, by design, never cues.
\end{itemize}

\begin{takeaway}
Read the rest of this report through one lens: Cadence is a \emph{faithful witness} to
freezing, not a treatment for it. Every choice that follows --- sensor, maths, detector,
display --- serves the single goal of detecting freezes honestly and showing them
clearly.
\end{takeaway}
